\begin{theorem}
Let $Q$ be a type and let $r$ be a quasi-order on $Q$.  Write
$\mathsf{HerCtblPower}(Q)$ for the hereditary countable non-empty part of
the full tagged hierarchy, and equip it with the induced relation
$\mathsf{HerCtblRel}(r)$.  Then $Q$ is better-quasi-ordered if and only if
$\mathsf{HerCtblPower}(Q)$ is well-founded for the strict-part relation
\[
(a,b)\longmapsto
\mathsf{HerCtblRel}(r)(a,b)\mathbin{\wedge}
\neg\mathsf{HerCtblRel}(r)(b,a).
\]
Here the strict part is oriented from the later, strictly smaller element
$a$ to the earlier element $b$, as in a descending chain.
\end{theorem}
\begin{proof}
Assume first that $Q$ is better-quasi-ordered.  By
\noderef{BqoHerCtblWqo}, the hereditary carrier equipped with
$\mathsf{HerCtblRel}(r)$ is well-quasi-ordered.  Applying
\noderef{HerCtblWqoIffWellFounded} to this quasi-order gives well-foundedness
of the relation
\[
(a,b)\longmapsto
\mathsf{HerCtblRel}(r)(a,b)\mathbin{\wedge}
\neg\mathsf{HerCtblRel}(r)(b,a),
\]
where a successive pair in a descending chain is therefore
$(a_{n+1},a_n)$.

Conversely, suppose that this strict-part relation on
$\mathsf{HerCtblPower}(Q)$ is well-founded.  The same equivalence
\noderef{HerCtblWqoIffWellFounded}, used in the reverse direction, implies
that $\mathsf{HerCtblRel}(r)$ is well-quasi-ordered.  The reverse direction
of \noderef{BqoHerCtblWqo} then implies that $Q$ is better-quasi-ordered.
Thus the two conditions are equivalent.
\end{proof}
